% =========================================================
%  ppt/02-game-loop-event-clock.tex
%  Materi 2 — Game Loop, Event Handling, dan Clock
%  Bab IV buku ajar
% =========================================================
\documentclass[aspectratio=169,10pt]{beamer}
% =========================================================
%  ppt/config/packages.tex
%  Paket bersama untuk seluruh presentasi kuliah
%  Pengantar Pemrograman Game
% =========================================================

\usepackage[T1]{fontenc}
\usepackage[utf8]{inputenc}
\usepackage[indonesian]{babel}
\usepackage{lmodern}
\usepackage{textcomp}
\usepackage{amsmath}
\usepackage{amssymb}
\usepackage{graphicx}
\usepackage{booktabs}
\usepackage{tabularx}
\usepackage{array}
\usepackage{multicol}
\usepackage{listings}
\usepackage{upquote}
\usepackage{adjustbox}
\usepackage{tcolorbox}
\usepackage{tikz}
\usetikzlibrary{
  arrows.meta,
  positioning,
  shapes.geometric,
  shapes.misc,
  calc,
  fit,
  backgrounds,
  matrix,
  chains,
  decorations.pathreplacing,
  patterns
}

\tcbuselibrary{skins}

% Jangan memuat geometry: ukuran halaman diatur oleh beamer.

% =========================================================
%  ppt/config/theme.tex
%  Satu tema Beamer untuk keempat belas presentasi
%  Pengantar Pemrograman Game
% =========================================================

% --------------------------------------------------------- warna
\definecolor{ubNavy}{HTML}{1F3864}
\definecolor{ubAccent}{HTML}{2E75B6}
\definecolor{ubHead}{HTML}{D9E2F3}
\definecolor{ubSoft}{HTML}{F2F6FC}
\definecolor{ubLine}{HTML}{BFBFBF}
\definecolor{ubText}{HTML}{1F2937}
\definecolor{ubMuted}{HTML}{5A6678}
\definecolor{ubGreen}{HTML}{1E7B45}
\definecolor{ubGreenBg}{HTML}{F1F8F2}
\definecolor{ubRed}{HTML}{B3261E}
\definecolor{ubRedBg}{HTML}{FDF2F1}
\definecolor{ubWarn}{HTML}{A35A05}
\definecolor{ubWarnBg}{HTML}{FFF8EC}
\definecolor{ubGray}{HTML}{F6F6F6}
\definecolor{ubCodeBg}{HTML}{10233F}
\definecolor{ubCodeText}{HTML}{E8EEF7}
\definecolor{ubCodeKey}{HTML}{8FBCF0}
\definecolor{ubCodeCom}{HTML}{8FB37A}
\definecolor{ubCodeStr}{HTML}{D9A46B}
\definecolor{ubCodeNum}{HTML}{7E99BE}

% Nama warna yang dipakai gambar TikZ buku ajar.
\definecolor{ubPrimary}{HTML}{1F3864}
\definecolor{ubAlt}{HTML}{F2F6FC}

% --------------------------------------------------------- dasar
\usetheme{default}
\usefonttheme{professionalfonts}
\setbeamertemplate{navigation symbols}{}
\setbeamerfont{frametitle}{size=\Large,series=\bfseries}
\setbeamerfont{framesubtitle}{size=\small}
\setbeamerfont{footline}{size=\tiny}
\setbeamerfont{block title}{size=\small,series=\bfseries}
\setbeamerfont{block body}{size=\small}
\setbeamerfont{itemize/enumerate body}{size=\small}
\setbeamerfont{itemize/enumerate subbody}{size=\footnotesize}

\setbeamersize{text margin left=0.60cm,text margin right=0.60cm}
\uchyph=0

\setbeamercolor{normal text}{fg=ubText,bg=white}
\setbeamercolor{structure}{fg=ubAccent}
\setbeamercolor{frametitle}{fg=ubNavy,bg=white}
\setbeamercolor{title}{fg=ubNavy}
\setbeamercolor{footline}{fg=ubMuted,bg=white}
\setbeamercolor{itemize item}{fg=ubAccent}
\setbeamercolor{itemize subitem}{fg=ubAccent}
\setbeamercolor{block title}{fg=white,bg=ubNavy}
\setbeamercolor{block body}{fg=ubText,bg=ubSoft}
\setbeamercolor{block title alerted}{fg=white,bg=ubWarn}
\setbeamercolor{block body alerted}{fg=ubText,bg=ubWarnBg}
\setbeamercolor{block title example}{fg=white,bg=ubGreen}
\setbeamercolor{block body example}{fg=ubText,bg=ubGreenBg}
\setbeamercolor{caption name}{fg=ubMuted}

\setbeamertemplate{blocks}[rounded][shadow=false]
\setbeamertemplate{itemize item}{\raisebox{0.10em}{\color{ubAccent}\scriptsize$\blacksquare$}}
\setbeamertemplate{itemize subitem}{\raisebox{0.06em}{\color{ubAccent}\tiny$\blacktriangleright$}}
\setbeamertemplate{enumerate item}{%
  \makebox[1.3em][l]{\color{ubAccent}\insertenumlabel.}}

\addtobeamertemplate{itemize/enumerate body begin}{%
  \setlength{\itemsep}{0.42em}\setlength{\parskip}{0pt}}{}

% --------------------------------------------------------- judul frame
\setbeamertemplate{frametitle}[default][left]

% --------------------------------------------------------- kaki
\newcommand{\materino}{00}
\newcommand{\materifoot}{Pengantar Pemrograman Game \textperiodcentered\ Materi \materino}

\setbeamertemplate{footline}{%
  \begin{beamercolorbox}[wd=\paperwidth,ht=0.015cm,dp=0cm]{footline}%
    \color{ubHead}\rule{\paperwidth}{0.4pt}%
  \end{beamercolorbox}%
  \begin{beamercolorbox}[wd=\paperwidth,ht=0.40cm,dp=0.14cm,%
                         leftskip=0.60cm,rightskip=0.60cm]{footline}%
    \usebeamerfont{footline}\color{ubMuted}\materifoot\hfill\insertframenumber%
  \end{beamercolorbox}%
}

% --------------------------------------------------------- makro buku ajar
\providecommand{\code}[1]{\texttt{\detokenize{#1}}}
\providecommand{\istilah}[1]{\emph{#1}}
\providecommand{\CPMK}[1]{CPMK\,#1}
\providecommand{\CPL}[1]{CPL\,#1}
\providecommand{\SUBCPMK}[1]{Sub-CPMK\,#1}

% --------------------------------------------------------- gaya kode
\lstdefinelanguage{CSharp}{
  morekeywords={abstract,as,base,bool,break,byte,case,catch,char,class,const,
    continue,decimal,default,delegate,do,double,else,enum,event,explicit,
    extern,false,finally,fixed,float,for,foreach,goto,if,implicit,in,int,
    interface,internal,is,lock,long,namespace,new,null,object,operator,out,
    override,params,private,protected,public,readonly,ref,return,sbyte,
    sealed,short,sizeof,stackalloc,static,string,struct,switch,this,throw,
    true,try,typeof,uint,ulong,unchecked,unsafe,ushort,using,virtual,void,
    volatile,while,var,get,set},
  sensitive=true,
  morecomment=[l]{//},
  morecomment=[s]{/*}{*/},
  morestring=[b]",
}

\lstdefinestyle{kode}{
  basicstyle=\ttfamily\footnotesize\color{ubCodeText},
  keywordstyle=\color{ubCodeKey}\bfseries,
  commentstyle=\color{ubCodeCom}\itshape,
  stringstyle=\color{ubCodeStr},
  numberstyle=\tiny\color{ubCodeNum},
  numbers=left,
  numbersep=7pt,
  backgroundcolor=\color{ubCodeBg},
  frame=none,
  showstringspaces=false,
  breaklines=false,
  tabsize=4,
  columns=fullflexible,
  keepspaces=true,
  xleftmargin=12pt,
  xrightmargin=2pt,
  aboveskip=0pt,
  belowskip=0pt,
  upquote=true,
}
\lstset{style=kode}

\newcommand{\kodejelas}[1]{%
  \vspace{0.16cm}%
  {\footnotesize\color{ubMuted}#1\par}%
}
\newcommand{\sumber}[1]{%
  \vspace{0.08cm}%
  {\scriptsize\color{ubMuted}\itshape #1\par}%
}
\newcommand{\awalan}[1]{%
  {\small\color{ubMuted}#1\par}\vspace{0.18cm}%
}

% --------------------------------------------------------- kotak
\newtcolorbox{kotak}[1][]{enhanced, colback=ubSoft, colframe=ubAccent,
  boxrule=0pt, leftrule=2.2pt, arc=1pt, left=6pt, right=6pt, top=4pt,
  bottom=4pt, fontupper=\footnotesize, #1}
\newtcolorbox{peringatan}[1][]{enhanced, colback=ubWarnBg, colframe=ubWarn,
  boxrule=0pt, leftrule=2.2pt, arc=1pt, left=6pt, right=6pt, top=4pt,
  bottom=4pt, fontupper=\footnotesize, #1}
\newtcolorbox{kiat}[1][]{enhanced, colback=ubGreenBg, colframe=ubGreen,
  boxrule=0pt, leftrule=2.2pt, arc=1pt, left=6pt, right=6pt, top=4pt,
  bottom=4pt, fontupper=\footnotesize, #1}

% --------------------------------------------------------- sampul
% \deckcover{nomor}{minggu}{judul}{sub-cpmk}{teknologi}{sumber}{blok}
\newcommand{\deckcover}[7]{%
\begin{frame}[plain]
\begin{tikzpicture}[remember picture, overlay]
  \fill[ubNavy] (current page.south west) rectangle ([xshift=6.00cm]current page.north west);
  \fill[ubAccent] ([xshift=6.00cm]current page.south west) rectangle ([xshift=6.08cm]current page.north west);

  \node[anchor=north west, text width=4.72cm, align=left, inner sep=0pt]
    at ($(current page.north west)+(0.66cm,-0.60cm)$) {%
      {\color{white}\bfseries\fontsize{10.2}{12.2}\selectfont
       Universitas Bale Bandung}\\[4pt]
      {\color{ubHead}\fontsize{7.6}{9.4}\selectfont
       Fakultas Teknologi Informasi}\\[1pt]
      {\color{ubHead}\fontsize{7.6}{9.4}\selectfont
       Program Studi Teknik Informatika}\\[7pt]
      {\color{ubAccent}\rule{4.0cm}{0.5pt}}%
    };

  \node[anchor=north west, text width=4.72cm, align=left, inner sep=0pt]
    at ($(current page.north west)+(0.66cm,-2.72cm)$) {%
      {\color{ubAccent}\fontsize{6.4}{7.8}\selectfont MATA KULIAH}\\[3pt]
      {\color{white}\bfseries\fontsize{9.8}{12.0}\selectfont
       Pengantar Pemrograman Game}\\[3pt]
      {\color{ubHead}\fontsize{7.6}{9.4}\selectfont
       MWCA009 \textperiodcentered\ 3 SKS}\\[8pt]
      {\color{ubHead}\fontsize{7.6}{9.4}\selectfont
       Semester Ganjil 2026/2027}%
    };

  \node[anchor=south west, text width=4.72cm, align=left, inner sep=0pt]
    at ($(current page.south west)+(0.66cm,0.52cm)$) {%
      {\color{ubAccent}\fontsize{6.4}{7.8}\selectfont DOSEN PENGAMPU}\\[3pt]
      {\color{white}\fontsize{8.4}{10.4}\selectfont
       Cecep Suwanda, S.Si., M.Kom.}%
    };

  \node[anchor=north west, text width=8.95cm, align=left, inner sep=0pt]
    at ($(current page.north east)+(-9.40cm,-0.92cm)$) {%
      {\color{ubAccent}\bfseries\fontsize{8.2}{10.2}\selectfont
       MATERI #1\ \ \textperiodcentered\ \ MINGGU #2}\\[8pt]
      {\color{ubNavy}\bfseries\fontsize{19.4}{23.4}\selectfont
       \hyphenpenalty=10000 \exhyphenpenalty=10000 \raggedright #3}%
    };

  \node[anchor=north west, text width=8.95cm, align=left, inner sep=0pt]
    at ($(current page.north east)+(-9.40cm,-4.30cm)$) {%
      \renewcommand{\arraystretch}{1.30}%
      \begin{tabular}{@{}>{\color{ubMuted}\fontsize{7.3}{9.0}\selectfont}p{1.85cm}
                        >{\color{ubText}\fontsize{8.1}{10.0}\selectfont}p{6.7cm}@{}}
        Sub-CPMK      & #4 \\
        Teknologi     & #5 \\
        Sumber materi & #6 \\
        Blok          & #7 \\
      \end{tabular}%
    };
\end{tikzpicture}
\end{frame}%
}

% --------------------------------------------------------- pembatas bagian
\newcommand{\sectionframe}[3]{%
\begin{frame}[plain]
\begin{tikzpicture}[remember picture, overlay]
  \fill[ubNavy] (current page.south west) rectangle (current page.north east);
  \fill[ubAccent] (current page.north west) rectangle ([xshift=0.10cm]current page.south west);
  \node[anchor=west, inner sep=0pt]
    at ($(current page.west)+(1.00cm,1.45cm)$)
    {\color{ubNavy!52!white}\bfseries\fontsize{58}{58}\selectfont #1};
  \node[anchor=west, inner sep=0pt, text width=11.6cm, align=left]
    at ($(current page.west)+(1.05cm,-0.55cm)$) {%
      {\color{ubAccent}\bfseries\fontsize{8.0}{10.0}\selectfont BAGIAN}\\[3pt]
      {\color{white}\bfseries\fontsize{20.5}{24.5}\selectfont #2}\\[5pt]
      {\color{ubHead}\small #3}%
    };
  \node[anchor=south east, inner sep=0pt]
    at ($(current page.south east)+(-0.60cm,0.32cm)$)
    {\color{ubHead}\fontsize{7.4}{9}\selectfont #1 / 6
     \ \ \textperiodcentered\ \ \insertframenumber};
\end{tikzpicture}
\end{frame}%
}

% --------------------------------------------------------- peta pertemuan
\newcommand{\blokpeta}[6]{%
  \ifnum#2=#3
    \node[rounded corners=3pt, fill=ubNavy, draw=ubNavy, line width=0.5pt,
          text width=3.02cm, minimum height=1.62cm, align=left, inner sep=5pt]
      at (#1,0) {%
        {\color{ubAccent}\bfseries\fontsize{6.8}{8}\selectfont #4}\\[1pt]
        {\color{white}\bfseries\fontsize{7.6}{9}\selectfont #5}\\[2pt]
        {\color{ubHead}\fontsize{6.4}{7.6}\selectfont #6}};
  \else
    \node[rounded corners=3pt, fill=ubSoft, draw=ubHead, line width=0.5pt,
          text width=3.02cm, minimum height=1.62cm, align=left, inner sep=5pt]
      at (#1,0) {%
        {\color{ubAccent}\bfseries\fontsize{6.8}{8}\selectfont #4}\\[1pt]
        {\color{ubNavy}\bfseries\fontsize{7.6}{9}\selectfont #5}\\[2pt]
        {\color{ubMuted}\fontsize{6.4}{7.6}\selectfont #6}};
  \fi
}

% \petapertemuan{blok aktif}{sebelumnya}{sesudahnya}{keluaran}
% \petapertemuan{blok aktif}{sebelumnya}{pertemuan ini}{sesudahnya}{keluaran}
\newcommand{\petapertemuan}[5]{%
\begin{frame}{Peta Pertemuan}
{\scriptsize\color{ubMuted}Kedudukan materi dalam empat blok perkuliahan}\par
\vspace{0.10cm}
\centering
\begin{tikzpicture}
  \blokpeta{-5.58}{1}{#1}{BLOK I}{Materi 1--5}{Game loop, vektor,\\fisika, collision}
  \blokpeta{-1.86}{2}{#1}{BLOK II}{Materi 6--9}{AI NPC: chase, FSM,\\pathfinding A*}
  \blokpeta{1.86}{3}{#1}{BLOK III}{Materi 10--11}{Unity: scene,\\level, physics}
  \blokpeta{5.58}{4}{#1}{BLOK IV}{Materi 12--14}{Aset visual, audio,\\UI, playtest}
\end{tikzpicture}
\par\vspace{0.30cm}
\begin{columns}[T,onlytextwidth]
  \begin{column}{0.315\textwidth}
    \begin{beamercolorbox}[wd=\linewidth,sep=4pt,rounded=true]{block body}
      {\color{ubMuted}\fontsize{6.6}{8}\selectfont SEBELUMNYA}\\[2pt]
      {\footnotesize #2}
    \end{beamercolorbox}
  \end{column}
  \begin{column}{0.315\textwidth}
    \begin{beamercolorbox}[wd=\linewidth,sep=4pt,rounded=true]{block title}
      {\color{ubHead}\fontsize{6.6}{8}\selectfont PERTEMUAN INI}\\[2pt]
      {\footnotesize #3}
    \end{beamercolorbox}
  \end{column}
  \begin{column}{0.315\textwidth}
    \begin{beamercolorbox}[wd=\linewidth,sep=4pt,rounded=true]{block body}
      {\color{ubMuted}\fontsize{6.6}{8}\selectfont SESUDAHNYA}\\[2pt]
      {\footnotesize #4}
    \end{beamercolorbox}
  \end{column}
\end{columns}
\par\vspace{0.26cm}
{\footnotesize\color{ubText}\textbf{Keluaran pertemuan.} #5\par}
\end{frame}%
}

% --------------------------------------------------------- penutup
\newcommand{\penutup}[2]{%
\begin{frame}[plain]
\begin{tikzpicture}[remember picture, overlay]
  \fill[ubNavy] (current page.south west) rectangle (current page.north east);
  \fill[ubAccent] (current page.north west) rectangle ([xshift=0.10cm]current page.south west);

  \node[anchor=north west, text width=13.6cm, align=left, inner sep=0pt]
    at ($(current page.north west)+(1.00cm,-0.95cm)$) {%
      {\color{ubAccent}\bfseries\fontsize{8.2}{10}\selectfont ASESMEN MINGGU INI}\\[6pt]
      {\color{white}\bfseries\fontsize{24}{28}\selectfont #1}%
    };

  \node[anchor=north west, text width=13.6cm, align=left, inner sep=0pt]
    at ($(current page.north west)+(1.00cm,-4.55cm)$) {%
      {\color{ubNavy!0!ubHead}\rule{13.4cm}{0.4pt}}\\[7pt]
      {\color{ubAccent}\bfseries\fontsize{7.4}{9}\selectfont
       PENGAIT KE MATERI BERIKUTNYA}\\[4pt]
      {\color{ubHead}\fontsize{10.6}{13.2}\selectfont #2}%
    };

  \node[anchor=south east, inner sep=0pt]
    at ($(current page.south east)+(-0.60cm,0.32cm)$)
    {\color{ubHead}\fontsize{7.4}{9}\selectfont \insertframenumber};
\end{tikzpicture}
\end{frame}%
}



\renewcommand{\materino}{02}

\begin{document}

\deckcover{02}{2}
  {Game Loop, Event Handling, dan Clock}
  {1.1}
  {Python / Pygame}
  {Bab IV buku ajar}
  {Blok I \textperiodcentered\ Materi 1--5}

% --------------------------------------------------------- peta
\petapertemuan{1}
  {Materi 1 \textperiodcentered\ Pengantar dan persiapan lingkungan}
  {Materi 2 \textperiodcentered\ Game Loop, Event Handling, dan Clock}
  {Materi 3 \textperiodcentered\ Gerak Objek, Vector2, Input, dan Delta Time}
  {Siklus permainan berjalan dengan laju bingkai terkendali, antrean
   peristiwa dikosongkan setiap bingkai, dan laju terukur tercatat pada
   lembar praktikum.}

% --------------------------------------------------------- capaian
\begin{frame}{Capaian Pertemuan}
\setbeamerfont{itemize/enumerate body}{size=\scriptsize}
\awalan{Sub-CPMK 1.1 beserta tiga indikatornya, dan tujuan Bab IV buku ajar.}
\begin{columns}[T,onlytextwidth]
  \begin{column}{0.50\textwidth}
    \begin{beamercolorbox}[wd=\linewidth,sep=5pt,rounded=true]{block title}
      {\fontsize{6.8}{8}\selectfont\color{ubHead}SUB-CPMK 1.1 \textperiodcentered\ C2}\\[2pt]
      {\footnotesize Menjelaskan struktur \emph{game loop} (input,
       \emph{update}, \emph{render}) dan peran manajemen waktu dalam
       permainan.}
    \end{beamercolorbox}
    \vspace{0.14cm}

    {\fontsize{7}{8.5}\selectfont\bfseries\color{ubAccent}INDIKATOR}\par
    \vspace{0.06cm}
    \begin{itemize}
      \item Menjelaskan urutan tahapan input--\emph{update}--\emph{render}
            dengan benar.
      \item Mengidentifikasi peran \texttt{pygame.time.Clock} dan pembatasan
            \emph{FPS}.
      \item Membedakan \emph{fixed timestep} dan \emph{variable timestep}
            secara konseptual.
    \end{itemize}
  \end{column}
  \begin{column}{0.46\textwidth}
    {\fontsize{7}{8.5}\selectfont\bfseries\color{ubAccent}TUJUAN BAB IV}\par
    \vspace{0.06cm}
    \begin{itemize}
      \item Menjelaskan ketiga tahap siklus permainan beserta urutan dan
            alasan urutannya.
      \item Membaca masukan pemain sebagai peristiwa tanpa menghentikan
            jalannya siklus.
      \item Menjelaskan peran \texttt{pygame.time.Clock} dalam mengukur selang
            waktu dan menahan laju bingkai.
      \item Menulis siklus lengkap yang berjalan sampai jendela ditutup,
            dengan laju bingkai yang ditentukan sendiri.
      \item Mengukur pengaruh batas laju terhadap beban prosesor dan memilih
            batas yang wajar.
      \item Menjelaskan mengapa siklus ini belum cukup untuk menghasilkan
            gerak, dan apa yang ditambahkan pada Bab V.
    \end{itemize}
  \end{column}
\end{columns}
\end{frame}

% ========================================================= bagian 1
\section{Kerangka Siklus Permainan}
\sectionframe{01}{Kerangka Siklus Permainan}
  {Tiga tahap yang selalu sama dan alasan urutannya}

\begin{frame}{Tiga Tahap yang Selalu Sama}
\awalan{Siklus permainan terdiri atas tiga tahap yang dikerjakan berurutan
lalu diulang terus-menerus.}
\centering
\adjustbox{max width=0.90\linewidth,max height=0.52\textheight}{%
  % chapters/04/figures/siklus-game-loop.tikz
% Diagram siklus input -> update -> render untuk Bab IV.
% Dipanggil dari section-*.tex dengan \input{figures/siklus-game-loop.tikz}
\begin{tikzpicture}[
  node distance=20mm,
  tahap/.style={
    rectangle, rounded corners=2pt, draw=ubPrimary, thick, fill=ubAlt,
    text width=27mm, align=center, minimum height=12mm
  },
  panah/.style={-{Stealth[length=2.6mm]}, thick, draw=ubAccent}
]
\node[tahap] (input) {\small\bfseries Input\\[1pt]\footnotesize baca peristiwa};
\node[tahap, right=of input] (update) {\small\bfseries Update\\[1pt]\footnotesize perbarui keadaan};
\node[tahap, right=of update] (render) {\small\bfseries Render\\[1pt]\footnotesize gambar layar};

\draw[panah] (input) -- (update);
\draw[panah] (update) -- (render);
\draw[panah] (render.south) -- ++(0,-7mm) -| (input.south);
\node[below=11mm of update, font=\footnotesize\itshape, text=ubPrimary]
  {diulang terus-menerus selama permainan berjalan};
\end{tikzpicture}
}
\vspace{0.10cm}
\begin{columns}[T,onlytextwidth]
  \begin{column}{0.32\textwidth}
    {\scriptsize\textbf{Masukan} --- membaca peristiwa sejak bingkai
     sebelumnya: tombol, tetikus, atau jendela yang hendak ditutup.}
  \end{column}
  \begin{column}{0.34\textwidth}
    {\scriptsize\textbf{Pembaruan} --- mengubah keadaan: memindahkan objek,
     memeriksa tabrakan, menambah skor, menentukan perilaku NPC.}
  \end{column}
  \begin{column}{0.30\textwidth}
    {\scriptsize\textbf{Penayangan} --- menggambar ulang seluruh jendela
     menurut keadaan yang baru diperbarui.}
  \end{column}
\end{columns}
\end{frame}

\begin{frame}{Mengapa Urutannya Demikian}
\awalan{Urutan tahapan bukan pilihan gaya penulisan, melainkan keharusan yang
lahir dari sifat pekerjaannya.}
\begin{itemize}
  \item \textbf{Masukan didahulukan} karena pembaruan memerlukan pengetahuan
        tentang kehendak pemain.
        \begin{itemize}
          \item Bila pembaruan lebih dahulu, tombol yang baru ditekan baru
                berpengaruh pada bingkai berikutnya; pemain merasakan tanggapan
                terlambat satu bingkai.
        \end{itemize}
  \item \textbf{Penayangan dikerjakan terakhir} karena pemain harus melihat
        keadaan akhir setelah seluruh perhitungan selesai.
        \begin{itemize}
          \item Menggambar di tengah perhitungan menampilkan keadaan setengah
                jadi.
        \end{itemize}
\end{itemize}
\begin{kotak}
Ketika jumlah objek bertambah, urutan yang tidak konsisten menimbulkan gejala
yang sulit dilacak: objek tampak bergerak mendahului tabrakannya, atau skor
bertambah sebelum kejadiannya sungguh terjadi. Pada Unity kelak, ketiga tahap
itu tetap ada walaupun sebagian dikelola mesin.
\end{kotak}
\end{frame}

\begin{frame}{Kerangka Siklus dalam Bentuk Kode}
\awalan{Tanpa isi setiap tahapnya, kerangka siklus hanya memerlukan satu
perulangan dengan tiga hal penting.}
\begin{enumerate}
  \item \textbf{Perulangan dikendalikan sebuah variabel.} \texttt{running}
        diubah menjadi salah ketika pemain meminta keluar, sehingga siklus
        berakhir dengan rapi.
  \item \textbf{Variabel keadaan berada di luar perulangan.} Keadaan harus
        tetap hidup dari satu bingkai ke bingkai berikutnya.
  \item \textbf{Ketiga tahap berada dalam satu perulangan yang sama.} Bukan
        tiga perulangan terpisah yang berjalan tidak serempak.
\end{enumerate}
\begin{kotak}
Apabila variabel keadaan dibuat di dalam perulangan, isinya kembali ke nilai
awal setiap bingkai dan permainan kehilangan ingatannya. Perulangan tanpa cara
keluar hanya dapat dihentikan dengan mematikan program secara paksa.
\end{kotak}
\end{frame}

% ========================================================= bagian 2
\section{Membaca Masukan sebagai Peristiwa}
\sectionframe{02}{Membaca Masukan sebagai Peristiwa}
  {Antrean peristiwa, event.get, dan jenis peristiwa lazim}

\begin{frame}{Mengapa event.wait() Tidak Cukup}
\awalan{Program Bab III menunggu peristiwa dengan
\texttt{pygame.event.wait()}. Cara itu tidak dapat dipakai di dalam siklus.}
\begin{itemize}
  \item \emph{Event handling} adalah cara program membaca peristiwa pemain
        tanpa menghentikan jalannya siklus.
  \item \texttt{event.wait()} \textbf{menghentikan} program sampai ada
        peristiwa yang terjadi.
  \item Bila pemain diam, tidak ada peristiwa datang; tahap pembaruan dan
        penayangan tidak pernah dikerjakan sehingga permainan membeku.
  \item Permainan yang benar tetap berjalan walaupun pemain sedang diam,
        karena objek di dalamnya terus bergerak dan waktu terus berjalan.
\end{itemize}
\begin{peringatan}
Pygame mengumpulkan peristiwa ke dalam antrean. \texttt{event.get()} mengambil
seluruh isinya tanpa menunggu, sedangkan \texttt{event.wait()} menunggu.
Hanya yang pertama boleh dipakai di dalam siklus.
\end{peringatan}
\end{frame}

\begin{frame}[fragile]{Menangani Peristiwa di Dalam Siklus}
\awalan{Perulangan membaca seluruh peristiwa yang masuk sejak bingkai
sebelumnya, lalu menanganinya satu per satu.}
\lstinputlisting[language=Python,firstline=129,lastline=141]
  {../buku_ajar/code/python/ch04_game_loop.py}
\kodejelas{\texttt{event.get()} mengosongkan antrean. \texttt{QUIT} dan ESC
mengakhiri siklus, sedangkan panah atas dan bawah mengubah batas laju bingkai.}
\end{frame}

\begin{frame}{Jenis Peristiwa yang Lazim}
{\scriptsize
\renewcommand{\arraystretch}{1.12}
\begin{tabularx}{\textwidth}{@{}>{\raggedright\arraybackslash}p{3.0cm}
    >{\raggedright\arraybackslash}p{3.9cm}
    >{\raggedright\arraybackslash}X@{}}
\toprule
\textbf{Jenis peristiwa} & \textbf{Keterangan tambahan} & \textbf{Dipakai untuk} \\
\midrule
\texttt{QUIT} & tidak ada &
  Menutup jendela: pemain meminta program berakhir. \\
\texttt{KEYDOWN} & \texttt{event.key}, \texttt{event.mod} &
  Tombol ditekan; terjadi sekali untuk setiap penekanan. \\
\texttt{KEYUP} & \texttt{event.key} &
  Tombol dilepas; berguna untuk tindakan yang berakhir. \\
\texttt{MOUSEBUTTONDOWN} & \texttt{event.pos}, \texttt{event.button} &
  Tombol tetikus ditekan, beserta letak penunjuk. \\
\texttt{MOUSEMOTION} & \texttt{event.pos}, \texttt{event.rel} &
  Penunjuk tetikus bergeser. \\
\bottomrule
\end{tabularx}}
\vfill
\begin{kotak}
Gunakan peristiwa untuk tindakan yang terjadi sekali (melompat, menembak,
membuka menu, keluar). Untuk tindakan yang berlangsung selama tombol ditahan
(berjalan, mengarahkan), periksa keadaan tombol dengan
\texttt{pygame.key.get\_pressed()}.
\end{kotak}
\end{frame}

% ========================================================= bagian 3
\section{Mengendalikan Laju dengan Clock}
\sectionframe{03}{Mengendalikan Laju dengan Clock}
  {Mengukur selang waktu dan menahan laju bingkai}

\begin{frame}{Dua Pekerjaan clock.tick()}
\awalan{Satu pemanggilan \texttt{clock.tick()} mengerjakan dua hal sekaligus.}
\begin{enumerate}
  \item \textbf{Mengukur selang waktu} yang benar-benar berlangsung sejak
        pemanggilan sebelumnya, dan mengembalikannya dalam milidetik. Nilai
        inilah yang kelak dipakai sebagai \emph{delta time}.
  \item \textbf{Menahan laju} dengan menunggu sampai selang waktu satu bingkai
        terpenuhi, apabila batas laju diberikan sebagai argumen.
\end{enumerate}
\[
  \mathrm{dt} = \frac{\mathrm{tick}(60)}{1000}\ \text{detik}
\]
\begin{kotak}
Karena mengukur sekaligus menahan, \texttt{tick()} diletakkan pada akhir
setiap putaran, yaitu setelah penayangan. Selang yang diukur harus mencakup
seluruh pekerjaan satu bingkai, bukan sebagiannya.
\end{kotak}
\end{frame}

\begin{frame}{Mengukur Laju yang Benar-benar Tercapai}
\awalan{Batas laju menyatakan keinginan, bukan kenyataan.}
\begin{itemize}
  \item \texttt{clock.get\_fps()} melaporkan laju yang benar-benar tercapai,
        dihitung dari rata-rata selang beberapa bingkai terakhir.
  \item Nilainya dipakai untuk menampilkan keadaan program kepada pemain dan
        menolong penulis program.
  \item Selisih besar antara batas laju dan laju terukur merupakan tanda
        paling awal bahwa program terlalu berat.
  \item Pada program bab ini, nilai tersebut ditampilkan pada judul jendela
        setiap bingkai.
\end{itemize}
\begin{kotak}
Judul jendela dipilih karena selalu terlihat, tidak memakan ruang gambar
permainan, dan dapat dibaca walaupun jendela sedang tidak difokuskan.
\end{kotak}
\end{frame}

\begin{frame}{Mengapa Laju Bingkai Perlu Dibatasi}
\awalan{Bukankah permainan yang lebih cepat lebih baik? Jawabannya tidak, dan
alasannya ada tiga.}
\begin{enumerate}
  \item \textbf{Perilaku permainan tidak seragam.} Bila gerak bergantung pada
        jumlah bingkai, permainan lebih sulit pada komputer yang lebih cepat.
  \item \textbf{Beban prosesor dan panas berlebihan.} Prosesor yang dipakai
        penuh tanpa manfaat menaikkan suhu dan mempercepat pemakaian baterai.
  \item \textbf{Tidak ada manfaat yang sepadan.} Mata manusia tidak dapat
        membedakan ribuan bingkai per detik, dan layar hanya menampilkan
        sejumlah gambar terbatas.
\end{enumerate}
\begin{kotak}
Siklus yang dibatasi 60 bingkai per detik memakai sekitar tiga persen satu inti
prosesor; siklus yang sama tanpa pembatasan memakai hampir seratus persen untuk
pekerjaan yang tidak terlihat pemain.
\end{kotak}
\begin{kiat}
Nilai batas yang wajar untuk pembelajaran adalah 60 bingkai per detik.
Turunkan menjadi 30 apabila program mulai berat.
\end{kiat}
\end{frame}

% ========================================================= bagian 4
\section{Siklus Lengkap}
\sectionframe{04}{Siklus Lengkap}
  {Menyatukan masukan, pembaruan, penayangan, dan pengendali waktu}

\begin{frame}{Susunan Siklus Lengkap}
\awalan{Siklus lengkap menambahkan pengendali waktu di akhir setiap putaran,
setelah ketiga tahap pekerjaan.}
\centering
\adjustbox{max width=0.90\linewidth,max height=0.52\textheight}{%
  % chapters/04/figures/siklus-dengan-pengendali-waktu.tikz
% Siklus lengkap: tiga tahap pekerjaan ditambah pengendali waktu yang
% mengukur selang waktu dan menahan laju bingkai.
\begin{tikzpicture}[
  tahap/.style={
    rectangle, rounded corners=2pt, draw=ubPrimary, thick, fill=ubAlt,
    align=center, text width=25mm, inner sep=3.5pt, minimum height=18mm,
    font=\footnotesize
  },
  waktu/.style={tahap, draw=ubAccent, fill=white},
  panah/.style={-{Stealth[length=2.4mm]}, thick, draw=ubAccent}
]

\node[tahap] (input) at (0,0)
  {\textbf{1. Masukan}\\[2pt]\scriptsize peristiwa dari\\pemain};
\node[tahap, right=7mm of input] (update)
  {\textbf{2. Pembaruan}\\[2pt]\scriptsize keadaan\\permainan};
\node[tahap, right=7mm of update] (render)
  {\textbf{3. Penayangan}\\[2pt]\scriptsize gambar ulang\\seluruh jendela};
\node[waktu, right=7mm of render] (clock)
  {\textbf{4. Pengendali waktu}\\[2pt]\scriptsize ukur selang,\\tahan laju bingkai};

\draw[panah] (input) -- (update);
\draw[panah] (update) -- (render);
\draw[panah] (render) -- (clock);
\draw[panah] (clock.south) -- ++(0,-8mm) -| (input.south);

\node[font=\scriptsize\itshape, ubPrimary, anchor=north] at (4.8,-2.2)
  {diulang selama permainan berjalan};
\end{tikzpicture}
}
\par\vspace{0.10cm}
\sumber{Keempat unsur ini dipertahankan pada seluruh bab berikutnya; yang
berubah adalah isi tahap pembaruan.}
\end{frame}

\begin{frame}[fragile]{Tahap Penayangan dan Pengendali Waktu}
\lstinputlisting[language=Python,firstline=155,lastline=168]
  {../buku_ajar/code/python/ch04_game_loop.py}
\kodejelas{\texttt{draw\_window} menggambar ulang seluruh jendela dan
\texttt{flip} menayangkannya. \texttt{clock.tick(fps\_limit)} di akhir putaran
menahan sisa waktu, lalu nilainya dibagi 1000 menjadi detik.}
\end{frame}

\begin{frame}{Tiga Kebiasaan yang Perlu Dibangun}
\begin{enumerate}
  \item \textbf{Gambar seluruh jendela dari awal setiap bingkai.}
        \texttt{screen.fill()} menimpa isi lama; tanpanya objek yang bergerak
        meninggalkan jejak di sepanjang lintasannya.
  \item \textbf{Simpan keadaan di luar perulangan.} \texttt{frame\_count},
        \texttt{fps\_limit}, dan \texttt{start\_ticks} harus hidup dari satu
        bingkai ke bingkai berikutnya.
  \item \textbf{Ubah keadaan hanya pada tahap pembaruan.} Fungsi penggambaran
        tidak mengubah apa pun, sehingga dapat dipanggil berapa kali pun tanpa
        mempengaruhi jalannya permainan.
\end{enumerate}
\begin{peringatan}
Kesalahan yang paling sering terjadi: menambahkan penghitung bingkai di dalam
fungsi penggambaran. Akibatnya penghitung bertambah setiap kali gambar
disusun, bukan setiap bingkai.
\end{peringatan}
\end{frame}

% ========================================================= bagian 5
\section{Mengukur dan Menyetel Laju}
\sectionframe{05}{Mengukur dan Menyetel Laju}
  {Membuktikan penghematan prosesor dengan pengukuran}

\begin{frame}[fragile]{Program Pengukuran}
\awalan{Siklus yang sama dijalankan tanpa jendela selama jangka waktu tetap;
setiap batas laju diuji bergantian.}
\lstinputlisting[language=Python,firstline=55,lastline=68]
  {../buku_ajar/code/python/ch04_ukur_laju.py}
\kodejelas{\texttt{time.perf\_counter()} mengukur waktu berjalan (jam dinding),
sedangkan \texttt{time.process\_time()} mengukur waktu prosesor. Perbandingan
keduanya menghasilkan pemakaian prosesor.}
\end{frame}

\begin{frame}{Hasil Pengukuran}
{\scriptsize
\renewcommand{\arraystretch}{1.14}
\begin{tabularx}{\textwidth}{@{}>{\centering\arraybackslash}p{1.7cm}
    >{\centering\arraybackslash}p{1.9cm}
    >{\centering\arraybackslash}p{2.3cm}
    >{\centering\arraybackslash}p{2.6cm}
    >{\raggedright\arraybackslash}X@{}}
\toprule
\textbf{Batas laju} & \textbf{Bingkai} & \textbf{Laju tercapai} &
\textbf{Waktu prosesor} & \textbf{Pemakaian prosesor} \\
\midrule
30    & 60      & 29,7 FPS    & 0,02 detik & 0,8\% \\
60    & 120     & 59,9 FPS    & 0,06 detik & 3,1\% \\
120   & 233     & 116,0 FPS   & 0,09 detik & 4,7\% \\
tanpa & 37.884  & 18.941 FPS  & 1,94 detik & 96,9\% \\
\bottomrule
\end{tabularx}}
\vspace{0.10cm}
\sumber{Pengukuran penulis pada lingkungan tanpa jendela, lama pengujian dua
detik setiap konfigurasi. Angka pada komputer Anda akan berbeda; yang penting
adalah hubungan antarbarisnya.}
\begin{kotak}
Batas laju dipatuhi dengan tepat, pemakaian prosesor tetap kecil pada semua
batas yang wajar, dan tanpa pembatasan prosesor dipakai penuh hampir
sepanjang waktu.
\end{kotak}
\end{frame}

\begin{frame}{Anggaran Waktu dan Pilihan Batas}
\awalan{Pekerjaan satu bingkai pada contoh: masukan 0,5 ms, pembaruan 5 ms,
dan penayangan 25 ms, sedangkan jatah pada batas 60 adalah}
\[
  \frac{1000\ \text{ms}}{60} = 16{,}7\ \text{ms per bingkai}.
\]
\begin{itemize}
  \item Jumlah pekerjaan 30,5 ms melebihi jatah; \texttt{clock.tick()} tidak
        menahan apa pun dan laju hanya mencapai sekitar 32,8 bingkai per detik.
  \item Menaikkan batas tidak mempercepat program yang pekerjaannya memang
        berat; yang menentukan adalah lama pekerjaan satu bingkai.
  \item \textbf{60 bingkai per detik} pilihan bawaan yang baik; \textbf{30}
        bila program berat; \textbf{120 ke atas} hanya bila benar-benar perlu.
\end{itemize}
\begin{kiat}
Batas laju bukan pembatasan mutu, melainkan penghematan yang tidak mengurangi
apa pun yang dapat dirasakan pemain. Ukur, jangan menebak.
\end{kiat}
\end{frame}

% ========================================================= bagian 6
\section{Yang Belum Ada pada Siklus Ini}
\sectionframe{06}{Yang Belum Ada pada Siklus Ini}
  {Dari siklus yang berjalan menuju objek yang bergerak}

\begin{frame}{Siklus Berjalan, tetapi Tidak Ada yang Bergerak}
\awalan{Siklus permainan sudah berjalan, tetapi permainannya belum.}
\begin{itemize}
  \item Program hanya memperbarui keadaan berupa penghitung bingkai; belum ada
        objek yang berpindah tempat.
  \item Perpindahan objek memerlukan \textbf{keadaan yang lebih kaya}: posisi
        dan kecepatan setiap objek harus disimpan dan diperbarui.
  \item Perpindahan juga memerlukan \textbf{pemakaian selang waktu}: nilai dari
        \texttt{clock.tick()} dipakai sebagai pengali kecepatan.
  \item Pemisahan ini disengaja agar kesalahan pada gerak nanti dapat
        ditelusuri tanpa mempertanyakan apakah siklusnya sudah benar.
\end{itemize}
\begin{kotak}
Perangkatnya sudah lengkap; yang belum ada hanyalah pemakaiannya. Bab V memakai
nilai selang waktu yang sama untuk menggerakkan objek, sehingga jarak yang
ditempuh sama setiap detik terlepas dari laju bingkai.
\end{kotak}
\end{frame}

\begin{frame}{Fixed Timestep dan Variable Timestep}
\begin{columns}[T,onlytextwidth]
  \begin{column}{0.48\textwidth}
    \begin{block}{Fixed timestep}
      \begin{itemize}
        \item Simulasi selalu maju dengan panjang langkah yang sama,
              $\Delta t = 1/60$ detik.
        \item Selisih waktu yang belum terpakai ditabung pada penampung.
        \item Hasilnya dapat diulang persis dan perhitungan fisika stabil.
      \end{itemize}
    \end{block}
  \end{column}
  \begin{column}{0.48\textwidth}
    \begin{block}{Variable timestep}
      \begin{itemize}
        \item Simulasi maju mengikuti lama bingkai sesungguhnya,
              $\Delta t = t_n - t_{n-1}$.
        \item Panjang langkahnya berubah-ubah; lebih sederhana diterapkan.
        \item Bila langkah menjadi sangat besar, objek cepat dapat melompat
              melewati rintangan tanpa terdeteksi.
      \end{itemize}
    \end{block}
  \end{column}
\end{columns}
\vspace{0.22cm}
\begin{kotak}
Siklus yang baru saja Anda tulis termasuk \emph{variable timestep}: setiap
putaran memakai selang yang benar-benar berlangsung, kadang 16 milidetik,
kadang 34 milidetik.
\end{kotak}
\end{frame}

\begin{frame}{Perkembangan Siklus pada Bab Berikutnya}
\awalan{Bab-bab selanjutnya menambahkan bagian pada siklus yang sama, bukan
menggantinya.}
{\scriptsize
\renewcommand{\arraystretch}{1.14}
\begin{tabularx}{\textwidth}{@{}>{\centering\arraybackslash}p{1.9cm}
    >{\raggedright\arraybackslash}p{5.2cm}
    >{\raggedright\arraybackslash}X@{}}
\toprule
\textbf{Bab} & \textbf{Yang ditambahkan} & \textbf{Akibat pada siklus} \\
\midrule
Bab V & Posisi, kecepatan, dan pemakaian selang waktu &
  Pembaruan mulai mengubah keadaan objek; penayangan menggambar posisi baru. \\
Bab VI & Gravitasi, lompatan, dan pantulan &
  Pembaruan menghitung percepatan, lalu kecepatan sebelum posisi. \\
Bab VII & \emph{Sprite}, \emph{Group}, dan tabrakan &
  Pembaruan memeriksa tabrakan dan mengubah keadaan permainan. \\
Bab VIII--XI & Pengendalian NPC &
  Pembaruan memanggil lapisan keputusan dan navigasi; sebagian tidak setiap
  bingkai. \\
Bab XIII & Fisika mesin pada Unity &
  Langkah tetap untuk fisika; kode Anda dipanggil mengikuti lama bingkai. \\
\bottomrule
\end{tabularx}}
\end{frame}

% --------------------------------------------------------- aktivitas
\begin{frame}{Aktivitas di Kelas \textperiodcentered\ Praktikum 02}
\awalan{Pengalaman belajar minggu 2 pada RPS, dikerjakan berurutan di kelas
lalu diselesaikan di rumah.}
\begin{enumerate}
  \item Tulis \texttt{ch04\_game\_loop.py} secara mandiri: siklus
        masukan--pembaruan--penayangan yang lengkap.
  \item Tangkap peristiwa keluar dan tombol: \texttt{QUIT}, ESC, panah
        atas/bawah, dan spasi.
  \item Batasi laju bingkai dengan \texttt{Clock} dan tampilkan laju terukur
        pada judul jendela.
  \item Ubah batas FPS dari 30 sampai 120, lalu catat laju terukur pada setiap
        nilai.
  \item Jalankan \texttt{ch04\_ukur\_laju.py} dan bandingkan pemakaian prosesor
        dengan dan tanpa pembatasan.
\end{enumerate}
\vspace{0.12cm}
{\footnotesize\color{ubAccent}\itshape Rincian lengkap: Praktikum 02 pada
Bab IV buku ajar. Asesmennya Tugas Praktikum 02 dengan observasi kinerja
\emph{game loop} (laju terukur dan pemakaian prosesor). Kumpulkan melalui
repositori Git pribadi.}
\end{frame}

% --------------------------------------------------------- rangkuman
\begin{frame}{Rangkuman}
\begin{itemize}
  \item Siklus permainan terdiri atas masukan, pembaruan, dan penayangan yang
        dikerjakan berurutan dan diulang selama permainan berjalan.
  \item Urutan tahapan bersifat wajib: masukan lebih dahulu agar pembaruan
        mengetahui kehendak pemain, penayangan paling akhir agar yang terlihat
        adalah keadaan yang sudah selesai dihitung.
  \item Variabel keadaan harus dibuat di luar perulangan agar nilainya bertahan
        dari satu bingkai ke bingkai berikutnya.
  \item \texttt{pygame.event.get()} mengambil seluruh isi antrean tanpa
        menunggu; \texttt{event.wait()} membekukan permainan, dan antrean harus
        dikosongkan setiap bingkai.
  \item \texttt{clock.tick()} mengukur selang waktu sekaligus menahan laju,
        sehingga diletakkan pada akhir putaran siklus.
  \item Batas laju adalah penghematan, bukan pembatasan mutu; menaikkan batas
        tidak mempercepat program yang pekerjaannya memang berat.
\end{itemize}
\end{frame}

% --------------------------------------------------------- cek
\begin{frame}{Cek Pemahaman}
\begin{enumerate}
  \item Jelaskan urutan ketiga tahap siklus permainan beserta pekerjaan pada
        masing-masing tahap.
  \item Apa dua pekerjaan yang dilakukan \texttt{clock.tick(60)}?
  \item Mengapa menaikkan batas laju tidak selalu menaikkan laju yang
        tercapai?
\end{enumerate}
\vspace{0.35cm}
{\footnotesize\color{ubMuted}\itshape Jawab tanpa membuka catatan, lalu
bandingkan dengan pasangan di sebelah Anda.}
\end{frame}

% --------------------------------------------------------- penutup
\penutup{Tugas Praktikum 02}
  {Materi 3 memakai selang waktu yang telah diukur untuk menggerakkan objek:
   posisi dan kecepatan sebagai \texttt{Vector2}, masukan tombol yang ditahan,
   dan \emph{delta time} agar kecepatan tetap konsisten pada 30 maupun 60 FPS.}

\end{document}
